This section validates the harness. With \harness{} in place, we
ran a 120-trial real-hardware study comparing the two
\rev{implementations} of the baseline: the static-scene LLM planner
and the camera-grounded VLM planner. The study answers two questions.
First, how reliable is the full three-layer system on real hardware?
Second, what fails once the five \rev{corrections} are active, and
does it match the one failure \harness{} does not \rev{resolve}?

\subsection{Design}
\label{sec:eval-design}

The study is anchored to the task and environment complexity taxonomy
and the TS/TSR/AETS metrics of Favali et al.~\cite{favaliRAS}, adapted
to a physical rig. Six task conditions cross three difficulty axes,
task length, reference specificity, and affordance and perception
load, each at an easy and a hard level, over the real scene of four
color-distinguished cubes and a tray:
\code{pp\_easy}/\code{pp\_hard} (pick-and-place, one vs.\ four
objects), \code{sort\_easy}/\code{sort\_hard} (stacking and sorting),
and \code{arith\_easy}/\code{arith\_hard} (the object to move must be
computed from a stated numeric mapping rather than named, isolating
reasoning load from perception load). Each condition ran 10 times per
\rev{implementation}, 120 trials in total. A single reasoning
strategy, chain-of-thought with self-consistency, was held fixed
throughout: the most safety-robust strategy in~\cite{favaliRAS},
chosen so that strategy variance is not stacked on top of the
comparison under test. Both \rev{implementations} use the same
underlying model (\code{kimi-k2.6}), text-only and image-grounded
respectively, so grounding modality is the only variable between
them.
The accompanying video presents the physical rig and
representative trials from all six conditions.

\subsection{Metrics, adapted for physical ground truth}
\label{sec:eval-metrics}

TS (task safety), TSR (task success rate), and AETS (action efficiency
on task success) are fully automatic in a symbolic simulator, where
damage and completion are exactly derivable from simulator state. Our
UR5cb has no collision or force-fault sensing in its control path, and
no vision-based verification of the final scene. We therefore adapt
the same three metrics to a semi-automatic protocol. Everything the
system can know automatically is logged per trial by the baseline's
debug instrumentation. Immediately after each trial, the operator
supplies exactly two numbers, whether the trial was safe and how many
sub-tasks completed, from which TS, TSR, and AETS are computed
identically to their original definitions.

At $n=10$ per cell, point estimates alone invite justified pushback,
so every reported TS and TSR carries a 95\% Wilson score interval (TS
is an exact per-trial proportion; TSR pools successes over
$\text{sub-tasks required} \times n$). One comparison was declared
before any data was collected: a two-sided Fisher's exact test, LLM
vs.\ VLM, on \code{arith\_hard} full-trial success, at $\alpha=0.05$.
It is binarized per trial rather than pooled across sub-tasks, since
sub-tasks are not independent: one wrong decision tends to misplace
every object in a trial together. Everything else is reported
descriptively.

\subsection{Results}
\label{sec:eval-results}

\begin{table}[t]
\centering
\caption{Overall results, 60 trials per \rev{planner implementation}
(120 total), with \harness{} in place.}
\label{tab:overall}
\begin{tabular}{@{}lcccc@{}}
\toprule
\textbf{\rev{Planner}} & \textbf{n} & \textbf{TS} & \textbf{TSR} & \textbf{AETS} \\
\midrule
LLM (\code{kimi-k2.6}) & 60 & 98.3\% & 98.3\% & 0.1269 \\
VLM (\code{kimi-k2.6}) & 60 & 85.0\% & 89.2\% & 0.1149 \\
\bottomrule
\end{tabular}
\end{table}

\begin{figure}[t]
\centering
\includegraphics[width=0.95\linewidth]{overall_summary.png}
\caption{Overall TS/TSR/AETS summary, LLM vs.\ VLM
\rev{implementation}, with \harness{} in place.}
\label{fig:overall}
\end{figure}

Table~\ref{tab:overall} reports the headline result. Both
\rev{implementations} share the same harness and the same underlying
model, so the residual gap between them is neither a grasp-geometry
gap nor a model-capability gap; it is attributable to grounding
modality. By task condition (Figure~\ref{fig:by-task}), the gap
concentrates on the two hardest multi-object conditions, and
predominantly in safety. TS is 100\% for both \rev{implementations} on
every easy condition and on \code{pp\_hard} and \code{sort\_hard} for
the LLM \rev{implementation}, but drops to 50\% (VLM,
\code{sort\_hard}) and 60\%/90\% (VLM/LLM, \code{arith\_hard}), the
conditions where the most objects share a zone at once. Neither
\rev{implementation} shows the sharp success-rate collapse an earlier
pilot on this task matrix, run with different per-modality models, had
found on the arithmetic conditions (30\%/25\% VLM TSR against
100\%/97\% LLM). With a single unified model, \code{arith\_easy} is
100\%/100\% and \code{arith\_hard} is 90\%/85\% (TSR, LLM/VLM): a real
but modest gap, not a collapse. The one pre-registered test, Fisher's
exact on \code{arith\_hard} full-trial success, shows no detectable
difference at this sample size (LLM 6/10, VLM 4/10, $p=0.66$),
consistent with a modest gap that $n=10$ per cell cannot resolve to
significance.

\begin{figure}[t]
\centering
\includegraphics[width=0.95\linewidth]{ts_tsr_by_task.png}
\caption{Task safety and task success rate by task condition, with
95\% Wilson intervals.}
\label{fig:by-task}
\end{figure}

\subsection{What fails with the harness in place}
\label{sec:eval-failures}

Reading all 22 non-empty operator failure notes from the 120 trials,
rather than trusting aggregates alone, was the key analysis step.
Every note describes either a collision between objects (10) or an
execution or grasp failure, a slipped grasp, a release too high, a
dropped object (12), during a multi-object task. Zero notes describe a
wrong object chosen for reasoning or value-mapping reasons, and zero
an incomplete enumeration (Figure~\ref{fig:failure-modes}). The
pattern holds in both \rev{implementations}: the LLM
\rev{implementation} contributes 4 of the 22 notes, all on
\code{arith\_hard}, and the VLM \rev{implementation} the remaining 18,
spread across the three multi-object conditions roughly in proportion
to trial counts.

\begin{figure}[t]
\centering
\includegraphics[width=0.85\linewidth]{failure_modes.png}
\caption{Failure-mode breakdown, derived from the 22 non-empty
operator notes out of 120 trials.}
\label{fig:failure-modes}
\end{figure}

This is the validation the section set out to obtain. With the
five corrections active, the dominant residual failure is
precisely the one \harness{} does not yet \rev{resolve}, the
sequential-placement collision, and it dominates \emph{both}
\rev{implementations}, regardless of which planner grounded the scene.
Zero notes implicate the model's choice of object or its task
decomposition. The remaining grasp and execution notes describe
physical instability during multi-object trials rather than the
\rev{resolved} height and matching classes.
% TODO-VERIFY: confermare dalle note operatore grezze che nessuna
% delle 12 note grasp/execution (es. "release too high") corrisponde
% in realtà a una classe di failure risolta; in tal caso riportarla
% come regressione parziale.
The residual failure is a property of the \rev{unresolved}
contact-dynamics problem, not of a perception modality.
